\section{Σύνοψη της εργασίας}
Η εργασία ξεκίνησε από ένα πρακτικό ερώτημα: πώς μπορεί μία πραγματική
εφαρμογή διαχείρισης κτηνιατρικής κλινικής να μεταφερθεί στο υπολογιστικό
νέφος χωρίς η υποδομή να γίνει δυσανάλογα ακριβή ή δύσκολη στη συντήρηση;
Το Happy Tails δεν αντιμετωπίστηκε ως απομονωμένο τεχνολογικό δείγμα, αλλά
ως σύστημα που χειρίζεται λογαριασμούς προσωπικού, κηδεμόνες, ζώα, ραντεβού,
ιατρικές εγγραφές και έγγραφα. Η αξιολόγηση έπρεπε επομένως να καλύπτει
λειτουργικότητα, ασφάλεια, διατήρηση δεδομένων, επιδόσεις, επαναφορά και κόστος.

Η τελική εγκατάσταση βασίστηκε σε μία EC2 instance στην περιοχή Frankfurt,
με Amazon Linux~2023 και τρία Docker containers: Nginx, Spring Boot και
MySQL. Το EBS παρείχε μόνιμη αποθήκευση, η Elastic IP και το Route~53
σταθερό δημόσιο σημείο πρόσβασης και το Nginx ασφαλή έκθεση μέσω HTTPS. Η
εφαρμογή συμπληρώθηκε με Amazon SES για τις ροές λογαριασμών προσωπικού και
με Google Calendar OAuth για τον συγχρονισμό ραντεβού. Οι δύο μηχανισμοί
παραμένουν διακριτοί: το SES μεταφέρει μηνύματα της εφαρμογής, ενώ το Google
Calendar δημιουργεί γεγονότα και προσκλήσεις.

Η εγκατάσταση δεν αξιολογήθηκε μόνο με στιγμιότυπα οθόνης. Ελέγχθηκαν οι
βασικές ροές, η επαναληψιμότητα των Liquibase migrations, η διατήρηση ενός
PDF μετά από επανεκκίνηση, οι ρόλοι, οι αποκρίσεις HTTP και η συμπεριφορά υπό
κλιμακούμενο φορτίο. Τα αποτελέσματα συσχετίστηκαν με CloudWatch metrics και
πραγματοποιήθηκε ελεγχόμενο reboot ολόκληρου του instance. Εξετάστηκαν ακόμη
οι εξαρτήσεις, το TLS και οι πραγματικοί παράγοντες κόστους. Το αποτέλεσμα
είναι συνεπώς τόσο μία λειτουργική εγκατάσταση όσο και μία αναπαραγώγιμη
διαδικασία αξιολόγησής της.

\section{Κύρια συμπεράσματα}
Το πρώτο συμπέρασμα είναι ότι μία μικρή containerized εγκατάσταση μπορεί να
καλύψει αξιόπιστα έναν περιορισμένο και σαφώς ορισμένο φόρτο. Δεν
παρατηρήθηκαν σφάλματα στα load stages, ακόμη και στη βαθμίδα των 25 VU. Η
υποβάθμιση της αρχικής αυθεντικοποίησης και η στιγμιαία αιχμή CPU έδειξαν
όμως ότι η χωρητικότητα δεν είναι απεριόριστη. Η καθυστέρηση μπορεί να γίνει
μη αποδεκτή πριν η εφαρμογή αρχίσει να επιστρέφει αποτυχημένες αποκρίσεις.

Το δεύτερο συμπέρασμα αφορά τη διαφορά ανάμεσα στην αυτόματη επανεκκίνηση
και στην υψηλή διαθεσιμότητα. Η πολιτική Docker
\texttt{restart: unless-stopped} επανέφερε τα services, το Liquibase δεν
επανεκτέλεσε ήδη εφαρμοσμένες αλλαγές και τα δεδομένα διατηρήθηκαν. Ωστόσο,
το reboot προκάλεσε διακοπή 40,79~s. Άρα υπάρχει μηχανισμός ανάκαμψης, αλλά
όχι πλεονασμός: όσο υπάρχει ένας μόνο κόμβος, η βλάβη ή συντήρησή του
επηρεάζει ολόκληρη την υπηρεσία.

Η επιλογή Spring Boot και React παρέμεινε κατάλληλη μετά τη μετάβαση στο
νέφος. Το Spring Boot παρέχει σαφή διαχωρισμό επιπέδων, ενσωμάτωση με Spring
Security και εξωτερική παραμετροποίηση~\cite{spring2026boot}. Η React
επιτρέπει επαναχρησιμοποιήσιμα στοιχεία διεπαφής και single-page εφαρμογή
\cite{react2026docs}. Ο συνδυασμός δεν εξαρτάται από μία συγκεκριμένη AWS
υπηρεσία: μπορεί να εκτελεστεί σε μία EC2 instance, σε πολλούς κόμβους ή ως
managed container workload. Αυτή η φορητότητα είναι σημαντικότερη από την
απλή χρήση πολλών cloud services.

Οι σαρώσεις έδειξαν επίσης ότι η συντήρηση των εξαρτήσεων είναι συνεχής
λειτουργική υποχρέωση. Εντοπίστηκαν διορθώσιμες ευπάθειες στο frontend tree
και στο deployed image. Η επόμενη έκδοση πρέπει επομένως να αναβαθμίσει το
Spring Boot dependency train, το React Router και το runtime base image και
να επαναλάβει τα regression, security και load tests.

Τέλος, το πραγματικό κόστος δεν ταυτίζεται με το ποσό που εμφανίζεται ως
πληρωτέο. Οι πιστώσεις οδήγησαν σε μηδενική πληρωμή, παρότι υπήρχε μετρήσιμη
χρήση. Η αξιολόγηση πρέπει να λαμβάνει υπόψη τις μονάδες χρέωσης EC2, EBS,
IPv4, Route~53, SES και CloudWatch, αλλά και τον χρόνο χειροκίνητης
διαχείρισης. Μία φθηνότερη υπηρεσία δεν είναι πάντοτε φθηνότερη συνολικά αν
μεταφέρει μεγάλο λειτουργικό βάρος στον διαχειριστή.

\section{Απάντηση στα ερευνητικά ερωτήματα}
\subsection*{ΕΕ1: Καταλληλότητα της οικονομικά περιορισμένης εγκατάστασης}
Η εγκατάσταση υποστήριξε τις βασικές λειτουργίες της κλινικής:
αυθεντικοποίηση και ρόλους, κηδεμόνες και ζώα, ραντεβού, ιατρικές εγγραφές,
έγγραφα, email λογαριασμού και ημερολόγιο. Η ανάπτυξη είναι επαναλήψιμη μέσω
Docker Compose και Liquibase και η δημόσια πρόσβαση προστατεύεται με HTTPS
και περιορισμένους κανόνες δικτύου. Η απάντηση στο ΕΕ1 είναι θετική για το
συγκεκριμένο εύρος και μέγεθος χρήσης, όχι όμως για απεριόριστο φορτίο ή
απαίτηση αδιάλειπτης λειτουργίας.

\subsection*{ΕΕ2: Επιδόσεις, πόροι και επανεκκίνηση}
Η εγκατάσταση διατήρησε μηδενικά σφάλματα στα load stages και ολοκλήρωσε το
soak test χωρίς απώλεια διαθεσιμότητας. Τα read endpoints παρέμειναν σχετικά
σταθερά, ενώ η αρχική αυθεντικοποίηση αποτέλεσε το πρώτο εμφανές σημείο
υποβάθμισης. Η CPU έφθασε στιγμιαία το 77,47\%, χωρίς παρατεταμένο κορεσμό.
Μετά το reboot η δημόσια υπηρεσία επανήλθε σε 40,79~s και τα δεδομένα
παρέμειναν αναλλοίωτα. Η συμπεριφορά ήταν σταθερή για το πειραματικό φορτίο,
αλλά ο χρόνος επαναφοράς επιβεβαιώνει τον περιορισμό του μοναδικού κόμβου.

\subsection*{ΕΕ3: Περιορισμοί και managed services}
Τα κύρια single points of failure είναι η μοναδική EC2 instance, η MySQL
στο ίδιο host και το τοπικά προσαρτημένο upload volume. Επιπλέον, τα μυστικά
βρίσκονται σε προστατευμένο αρχείο περιβάλλοντος, η παρακολούθηση βασίζεται
κυρίως σε προεπιλεγμένες μετρικές και το deployment περιλαμβάνει χειροκίνητα
βήματα. Οι περιορισμοί μπορούν να μειωθούν με Application Load Balancer
(ALB) και Auto Scaling σε περισσότερες Availability Zones, RDS Multi-AZ, S3 για τα
uploads, Secrets Manager, CloudWatch alarms, AWS Backup και WAF. Η επιλογή
πρέπει να γίνεται ανάλογα με την κρισιμότητα και όχι μόνο επειδή μία υπηρεσία
είναι διαθέσιμη.

\subsection*{ΕΕ4: Κόστος και συμβιβασμοί}
Η υλοποιημένη λύση έχει μικρό άμεσο κόστος και είναι εύκολη στην κατανόηση,
αλλά συγκεντρώνει compute, database και αρχεία στον ίδιο κόμβο. Μία λύση
υψηλότερης διαθεσιμότητας αυξάνει το σταθερό κόστος μέσω load balancer,
πολλαπλών targets, managed database, backups και αυξημένης παρακολούθησης.
Σε αντάλλαγμα μειώνει τον κίνδυνο διακοπής και επιτρέπει ανεξάρτητη κλιμάκωση
εφαρμογής και δεδομένων. Για μία μικρή κλινική η παρούσα λύση είναι λογικό
πρώτο στάδιο. Για πολλές κλινικές ή 24ωρη επιχειρησιακή εξάρτηση, το πρόσθετο
κόστος του πλεονασμού γίνεται δικαιολογημένο.

\section{Περιορισμοί}
Η μελέτη αφορά πραγματική αλλά μικρής κλίμακας εγκατάσταση. Το dataset, οι
λογαριασμοί και η διάρκεια των δοκιμών είναι περιορισμένα. Τα VU εκτελούν
συγκεκριμένες ακολουθίες HTTP και δεν προσομοιώνουν πλήρως την ανθρώπινη
συμπεριφορά, αργές συνδέσεις ή τη μακροχρόνια αύξηση δεδομένων. Ένα soak test
πέντε λεπτών μπορεί να εντοπίσει άμεσα προβλήματα, όχι όμως διαρροές πόρων που
εμφανίζονται μετά από ώρες ή ημέρες.

Το reboot αξιολόγησε την επαναφορά του υπάρχοντος instance, όχι την απώλεια
Availability Zone. Δεν πραγματοποιήθηκε restore σε νέο περιβάλλον ούτε
μετρήθηκαν Recovery Time Objective (RTO) και Recovery Point Objective (RPO)
μέσα από πλήρη άσκηση καταστροφής. Η επιτυχής persistence δεν αποτελεί από
μόνη της απόδειξη επαρκούς στρατηγικής backup.

Οι έλεγχοι ασφάλειας δεν αντικαθιστούν penetration test, ανεξάρτητη
ανασκόπηση κώδικα ή έλεγχο συμμόρφωσης. Για έγγραφα και στοιχεία επικοινωνίας
απαιτούνται πολιτικές διατήρησης, audit και ελαχιστοποίησης δεδομένων που
υπερβαίνουν το πεδίο της εργασίας. Τέλος, το κόστος επηρεάζεται από credits,
Free Tier, στρογγυλοποίηση και καθυστερημένη ενημέρωση billing. Αποτυπώνει
τη συγκεκριμένη περίοδο και όχι δεσμευτική μελλοντική πρόβλεψη.

\section{Μελλοντικές επεκτάσεις}
\subsection{Άμεσα βήματα με περιορισμένο πρόσθετο κόστος}
Πρώτη προτεραιότητα είναι η σκλήρυνση της υπάρχουσας λύσης: αναβάθμιση των
ευάλωτων dependencies, αυτοματοποιημένο build και deployment μέσω πρακτικών
Continuous Integration/Continuous Delivery (CI/CD), τακτικά EBS
snapshots και δοκιμασμένο restore. Τα μυστικά μπορούν να μεταφερθούν από το
\texttt{.env} στο AWS Secrets Manager, το οποίο υποστηρίζει ασφαλή αποθήκευση,
ανάκτηση και rotation credentials~\cite{aws2026secretsmanager}. CloudWatch
alarms για CPU, status checks, χώρο δίσκου και αποτυχίες εφαρμογής μπορούν να
ειδοποιούν πριν ένα συμβάν γίνει ορατό στους χρήστες
\cite{aws2026cloudwatchalarms}.

Τα uploads μπορούν να μεταφερθούν σε S3, ώστε να αποσυνδεθούν από τον κύκλο
ζωής του instance, με προσωρινά εξουσιοδοτημένα URLs και ιδιωτικό bucket. Η
έξοδος του SES από το sandbox πρέπει να συνοδευτεί από χειρισμό bounces και
complaints. Στο frontend μπορούν να προστεθούν καλύτερη προσβασιμότητα,
σαφείς καταστάσεις σφάλματος και end-to-end browser tests.

\subsection{Αρχιτεκτονική με μεγαλύτερο προϋπολογισμό}
Με μεγαλύτερο προϋπολογισμό, η σημαντικότερη αλλαγή είναι η αφαίρεση του
μοναδικού σημείου αστοχίας. Ένα Application Load Balancer μπορεί να
κατανέμει κίνηση μόνο σε υγιή targets και ένα Auto Scaling Group να προσθέτει
ή να αντικαθιστά instances σε περισσότερες Availability Zones
\cite{aws2026elbasg}. Το Spring Boot application layer πρέπει τότε να γίνει
stateless, δηλαδή να μη διατηρεί κρίσιμη κατάσταση σε συγκεκριμένο instance:
sessions, uploads και secrets δεν μπορούν να εξαρτώνται από τον τοπικό δίσκο
ενός container.

Η MySQL μπορεί να μεταφερθεί σε RDS Multi-AZ. Στη βασική μορφή διατηρείται
συγχρονισμένο standby σε άλλη ζώνη και παρέχεται failover χωρίς αλλαγή του
application endpoint~\cite{aws2026rdsmultiaz}. Αυτό δεν καταργεί την ανάγκη
backup, αλλά μειώνει σημαντικό λειτουργικό βάρος. Το AWS Backup μπορεί να
ενοποιήσει προγράμματα λήψης και διατήρησης αντιγράφων
\cite{aws2026backup}.

Για συχνότερες εκδόσεις, τα containers μπορούν να μεταφερθούν στο Amazon
Elastic Container Service (ECS). Το AWS Fargate εκτελεί τα tasks χωρίς άμεση
διαχείριση hosts, αλλά απαιτεί σαφείς task definitions,
όρια πόρων, κεντρικά logs και διαφορετικό μοντέλο κόστους
\cite{aws2026fargate}. Δεν είναι υποχρεωτικά φθηνότερο από μία μικρή EC2· έχει
αξία όταν η αυτοματοποίηση και η ανεξάρτητη κλιμάκωση δικαιολογούν το κόστος.

Στο δημόσιο άκρο, CloudFront και AWS WAF μπορούν να λειτουργήσουν μπροστά από
το ALB. Το WAF επιθεωρεί HTTP(S) αιτήματα και μπορεί να τα καταγράφει, να τα
απορρίπτει ή να εφαρμόζει challenge~\cite{aws2026waf}. Η ενεργοποίηση πρέπει
να αρχίσει σε count mode ώστε να μη διακοπούν νόμιμες ροές από υπερβολικά
γενικούς κανόνες.

\begin{table}[htbp]
\centering
\caption{Προτεινόμενη σταδιακή εξέλιξη της υποδομής.}
\label{tab:future-roadmap}
\small
\begin{tabularx}{\textwidth}{@{}p{2.7cm}p{4.2cm}X@{}}
\toprule
\textbf{Στάδιο} & \textbf{Κύρια αλλαγή} & \textbf{Αναμενόμενο όφελος} \\
\midrule
Άμεσο & Updates, CI/CD, alarms και δοκιμασμένο backup & Μείωση κινδύνου χωρίς αλλαγή αρχιτεκτονικής \\
Επόμενο & S3 uploads και Secrets Manager & Αποδέσμευση δεδομένων και μυστικών από τον host \\
Υψηλή διαθεσιμότητα & ALB, πολλαπλά targets και RDS Multi-AZ & Αφαίρεση των κύριων single points of failure \\
Μεγαλύτερη κλίμακα & ECS/Fargate, Auto Scaling, CloudFront και WAF & Αυτοματοποιημένη κλιμάκωση και προστασία του δημόσιου άκρου  \\
\bottomrule
\end{tabularx}
\end{table}

\subsection{Εξέλιξη του Spring Boot και της React εφαρμογής}
Στο Spring Boot μπορούν να προστεθούν application metrics, correlation IDs,
δομημένα logs και διακριτούς ελέγχους ετοιμότητας (readiness) και ζωτικότητας
(liveness). Ο πρώτος δηλώνει αν το instance μπορεί να δεχθεί κίνηση, ενώ ο
δεύτερος αν η διεργασία λειτουργεί και πρέπει να παραμείνει ενεργή. Οι εργασίες email
και Calendar μπορούν να εκτελούνται ασύγχρονα μέσω ουράς, ώστε προσωρινή
αποτυχία παρόχου να μην αποτυγχάνει τη συναλλαγή του χρήστη. Κλειδιά
ιδιοδυναμίας (idempotency keys), με τα οποία η επανάληψη του ίδιου αιτήματος
δεν δημιουργεί δεύτερο αποτέλεσμα, και ελεγχόμενη πολιτική επαναπροσπάθειας
θα μειώσουν τον κίνδυνο διπλών προσκλήσεων ή μηνυμάτων.

Στη React διεπαφή μπορούν να προστεθούν offline-friendly συμπεριφορά,
λεπτομερέστερες καταστάσεις δικαιωμάτων, προσβασιμότητα βάσει των Web Content
Accessibility Guidelines (WCAG) και
αυτοματοποιημένα component και end-to-end tests. Μελλοντική mobile ή
progressive web εφαρμογή έχει νόημα μόνο αν προκύπτει από τις πραγματικές
ανάγκες κτηνιάτρων και κηδεμόνων, όχι ως τεχνολογική προσθήκη χωρίς χρήση.

\subsection{Πού προσφέρει μεγαλύτερη αξία η AWS}
Η AWS υπερέχει όταν απαιτούνται ελαστικότητα, γεωγραφική διασπορά, managed
υπηρεσίες και αυτοματοποίηση λειτουργίας. Είναι ιδιαίτερα χρήσιμη όταν πολλές
κλινικές μοιράζονται την πλατφόρμα, υπάρχουν αιχμές κίνησης, απαιτούνται
αυστηρά RTO/RPO ή μία μικρή ομάδα δεν μπορεί να συντηρεί μόνη της database
replication, load balancers, backups και κεντρική παρακολούθηση. Το AWS
Well-Architected Framework βοηθά να εξετάζονται συστηματικά οι συμβιβασμοί
ασφάλειας, αξιοπιστίας, επίδοσης και κόστους~\cite{aws2024wellarchitected}.

Δεν υπερέχει όμως επειδή διαθέτει απλώς πολλές υπηρεσίες. Για μία μικρή
εφαρμογή με προβλέψιμο φορτίο, η μονοκομβική λύση μπορεί να είναι
καταλληλότερη, εφόσον υπάρχουν backup, monitoring και αποδεκτό σχέδιο
διακοπής. Η ώριμη αρχιτεκτονική δεν χρησιμοποιεί κάθε δυνατότητα του
παρόχου· αγοράζει με μετρήσιμο τρόπο την αξιοπιστία και την αυτοματοποίηση
που πράγματι χρειάζεται.

\subsection{Επόμενη πειραματική αξιολόγηση}
Μία επόμενη μελέτη θα μπορούσε να συγκρίνει την παρούσα εγκατάσταση με μία
Multi-AZ εκδοχή μέσω των ίδιων p50, p95, p99, throughput, error rate, CPU και
memory metrics. Χρειάζονται soak tests πολλών ωρών, failover της βάσης,
αφαίρεση application target υπό φορτίο και πλήρες restore από backup. Ο
συνδυασμός τεχνικών αποτελεσμάτων και μηνιαίου κόστους θα δείξει όχι μόνο
ποια λύση είναι ταχύτερη, αλλά πόσο κοστίζει κάθε ουσιαστική βελτίωση
διαθεσιμότητας.

Συνολικά, η εργασία δείχνει μία εφικτή διαδρομή: πρώτα κατανοητή και
μετρήσιμη υποδομή, έπειτα σκλήρυνση και αυτοματοποίηση και, όταν οι απαιτήσεις
το δικαιολογούν, πλεονασμός και managed orchestration. Έτσι το Happy Tails
μπορεί να αναπτυχθεί χωρίς να χάσει την απλότητα που επέτρεψε την αρχική του
λειτουργία.
